\section{AI Usage Documentation}

I used Anthropic's Claude throughout this project, working directly in the
repository through an agentic coding session with file and shell access.

For the code, I used it to scaffold the first version of the
\texttt{docker-compose.yml} layout, the controller's FastAPI skeleton and the
four generators' REST contract (start, stop, patch config, get status), and
then iterated on that by hand and through more prompts. The Adaptive Control
loop, the warmup/cooldown ramp, the random pattern on all four generators and
the per-protocol fault injection were built the same way, one function at a
time, with a Python syntax check after every change. The OpenAPI/Swagger
documentation was generated from the Pydantic models and FastAPI decorators
and kept in sync with a small regeneration script run after every API change.

For the documentation, I had it review the README and the docs against the
actual code and fix what no longer matched, for example an old claim that the
TCP/UDP generator used Scapy when the real implementation just uses plain
sockets.

This report itself was drafted the same way: the architecture, implementation
and configuration sections came from the verified source code, while the
traffic analysis, failure analysis and multi-machine sections were left as
placeholders on purpose, since those need real Wireshark captures that only I
can produce on lab hardware.

I reviewed every change before accepting it and ran the syntax checks myself.
I also went through the entire codebase module by module in a separate
walkthrough session to make sure I could explain every part of it on my own,
not just recognize that it works.

\todoblock{INSERT}{Replace this paragraph with your own honest account before
submission, including anything I missed above.}
